%Chargement des paquets
\usepackage{amsmath}
\usepackage{amsfonts}
\usepackage{amssymb}
\usepackage{enumerate}
\usepackage{mathtools}
\usepackage{bbm}
\usepackage{xparse, etoolbox}
\usepackage{enumerate}
\usepackage{enumitem}
\usepackage{mathabx}
\usepackage{babel}[french]

%environment
\usepackage{amsthm}
\usepackage{keytheorems}
\usepackage{hyperref} 

%Interface théorème
\newkeytheorem{lem}[
	name = Lemme
]

\newkeytheorem{prop}[
	name = Proposition
]

\newkeytheorem{thm}[
	name = Théorème
]

\newkeytheorem{coro}[
	name = Corollaire
]

\renewcommand*{\proofname}{Démonstration}

\theoremstyle{definition}
\newtheorem{defn}{Définition}

\theoremstyle{definition}
\newtheorem*{nota}{Notation}

\theoremstyle{definition}
\newtheorem*{limi}{Liminaire}

\theoremstyle{remark}
\newtheorem{rmq}{Remarque}

\theoremstyle{plain}
\newtheorem*{fait}{Fait}


%Ensembles de nombres
\newcommand{\N} {\mathbb{N}}
\newcommand{\Ne}{\N^\ast}
\newcommand{\Z} {\mathbb{Z}}
\newcommand{\Q} {\mathbb{Q}}
\newcommand{\R} {\mathbb{R}}
\newcommand{\Rb}{\overline{\mathbb{R}}}
\newcommand{\Rp}{\R_+}
\newcommand{\Rm}{\R_-}
\newcommand{\K} {\mathbb{K}}
\newcommand{\Ket} {\mathbb{K^{\ast}}}
\newcommand{\Cx}{\mathbb{C}}

%Ensembles, tribus
\newcommand{\A}{\mathbb{A}}
\renewcommand{\P}{\mathcal{P}}
\newcommand{\T}{\mathcal{T}}
\newcommand{\F}{\mathcal{F}}
\newcommand{\C} {\mathcal{C}}
\newcommand{\ouv}{\mathcal{O}}
\newcommand{\B}{\mathcal{B}}
\newcommand{\M}{\mathcal{M}}
\newcommand{\etag}{\mathcal{E}}
\newcommand{\etagOp}{\etag^{+}(\Omega)}
\newcommand{\etagO}{\etag(\Omega)}
%%Lp avant quotientage
\newcommand{\Lic}{\mathcal{L}^1}
\newcommand{\Lpc}{\mathcal{L}^p}
\newcommand{\Linfc}{\mathcal{L}^{\infty}}
\newcommand{\Lifull}{\Lic(\Omega, \B, m)}
\newcommand{\Lpfull}{\Lpc(\Omega, \B, m)}
\newcommand{\Linffull}{\Linfc(\Omega, \B, m)}
%%Lp après quotientage
\newcommand{\Li}{L^1}
\newcommand{\Ld}{L^2}
\newcommand{\Lp}{L^p}
\newcommand{\Linf}{L^{\infty}}
\newcommand{\Listd}{\Li(\Omega, m)} 
\newcommand{\Ldstd}{\Ld(\Omega, m)} 
\newcommand{\Lpstd}{\Lp(\Omega, m)} 

%Quelques opérateurs
\DeclareMathOperator{\codim}{codim}
\DeclareMathOperator{\rg}{rg}
\DeclareMathOperator{\tr}{tr}
\DeclareMathOperator{\vect}{Vect}
\DeclareMathOperator{\ima}{Im}
\DeclareMathOperator{\id}{Id}
\DeclareMathOperator{\sign}{sign}
\DeclareMathOperator{\ext}{ext}
\newcommand{\ind}{\mathbbm{1}}
\newcommand*{\spr}[2]{\left(\,#1\,\middle|\,#2\,\right)}
\newcommand{\interior}[1]{%
  {\kern0pt#1}^{\mathrm{o}}%
}
\DeclareMathOperator{\card}{card}
\DeclareMathOperator{\ppcm}{ppcm}

%Commandes de pure flemme
\newcommand{\veps}{\varepsilon}
\newcommand{\intab}{\int_{a}^{b}}
\newcommand{\intR}{\int_{-\infty}^{\infty}}
\newcommand{\intinf}{\int_{- \infty}^{+ \infty}}
\newcommand{\equi}{\Leftrightarrow}
\newcommand{\Kev}{$\K$-espace vectoriel }
\newcommand{\Kevs}{$\K$-espaces vectoriels }
\newcommand{\Rev}{$\R$-espace vectoriel }
\newcommand{\Revs}{$\R$-espaces vectoriels }
\newcommand{\Cev}{$\Cx$-espace vectoriel }
\newcommand{\Cevs}{$\Cx$-espaces vectoriels }
\newcommand{\evn}{(E, \norm{.})}
\newcommand{\interf}[2]{[#1,#2]}
\newcommand{\limb}{\lim\limits_}
\newcommand{\lebn}{\lambda_n}
\newcommand{\pp}{\text{~presque partout}}
\newcommand{\derivp}[2]{\frac{\partial #1}{\partial #2}}
\newcommand{\famiu}{(u_i)_{i \in I}}
\newcommand{\sev}{\text{~s.e.v.}}
\newcommand{\Mnp}{M_{n,p}(\K)}
\newcommand{\Mn}{M_{n}(\K)}
\newcommand{\tq}{\text{~t.q.~}}
\newcommand{\mq}{~montrer~que~}
\newcommand{\et}{\quad \text{et} \quad}
\newcommand{\ou}{\text{~ou~}}
\newcommand{\Kx}{\K[X]}


%Commandes de gars chelou
\renewcommand{\subset}{\subseteq}

%Commande restriction, ty duckduckgo
\def\restr#1#2{\mathchoice
              {\setbox1\hbox{${\displaystyle #1}_{\scriptstyle #2}$}
              \restraux{#1}{#2}}
              {\setbox1\hbox{${\textstyle #1}_{\scriptstyle #2}$}
              \restraux{#1}{#2}}
              {\setbox1\hbox{${\scriptstyle #1}_{\scriptscriptstyle #2}$}
              \restraux{#1}{#2}}
              {\setbox1\hbox{${\scriptscriptstyle #1}_{\scriptscriptstyle #2}$}
              \restraux{#1}{#2}}}
\def\restraux#1#2{{#1\,\smash{\vrule height .8\ht1 depth .85\dp1}}_{\,#2}} 

%Quelques normes
\DeclarePairedDelimiter\abs{\lvert}{\rvert}
\DeclarePairedDelimiter\labs{\bigl\lvert}{\bigr\rvert}
\DeclarePairedDelimiter\norm{\lVert}{\rVert}
\DeclarePairedDelimiterXPP\normi[1]{}\lVert\rVert{_1}{\ifblank{#1}{\:\cdot\:}{#1}}
\DeclarePairedDelimiterXPP\normd[1]{}\lVert\rVert{_2}{\ifblank{#1}{\:\cdot\:}{#1}}
\DeclarePairedDelimiterXPP\normp[1]{}\lVert\rVert{_p}{\ifblank{#1}{\:\cdot\:}{#1}}
\DeclarePairedDelimiterXPP\normq[1]{}\lVert\rVert{_q}{\ifblank{#1}{\:\cdot\:}{#1}}
\DeclarePairedDelimiterXPP\norminf[1]{}\lVert\rVert{_\infty}{\ifblank{#1}{\:\cdot\:}{#1}}

%Bracket en angle
\DeclarePairedDelimiter\angl{\langle}{\rangle}

%Set, parenthèses
\newcommand{\set}[1]{\{ #1 \}}
\newcommand{\bigset}[1]{\bigl\{ #1 \bigr\}}
\newcommand{\Bigset}[1]{\Bigl\{ #1 \Bigr\}}
\newcommand{\bigpar}[1]{\bigl( #1 \bigr)}
\newcommand{\Bigpar}[1]{\Bigl( #1 \Bigr)}

%Opitmisation des opérateurs de comparaison
\DeclarePairedDelimiter\tio{o (}{)}
\DeclarePairedDelimiter\gro{O (}{)}


\pagestyle{fancy}

\fancyhead[C]{}
\fancyhead[R]{}

\begin{document}

\underline{Source :} Gourdon - Analyse - page 41 puis 75 puis 327

\underline{Leçons :}
\begin{itemize}
\item 
\end{itemize}

\begin{thm}
Les espaces connexes de $\R$ sont les intervalles de $\R$.
\end{thm}

\begin{proof}
On montre dans un premier temps qu'un connexe $C$ de $\R$ est un intervalle. 
Par l'absurde, si $C$ n'est pas un intervalle on peut trouver $a,b \in C$ et $x \in \R$ tels que $a<x<b$ et $x \not \in C$.
On a alors $C \subset ]-\infty; x [ \cup ] x ; \infty [$ et donc $C$ n'est pas connexe. \\

Réciproquement, soit $I$ un intervalle de $\R$. $I$ est connexe s'il est réduit à un point, on suppose donc $a<b$.
On distingue plusieurs cas.
\begin{enumerate}[label=(\roman*)]
\item 
Si $I=]a,b[$ avec $-\infty \leq a < b \leq \infty$. 
On muni l'ensemble $\set{0,1}$ de la distance discrète $\delta$ et on considère l'application continue $f: I \to \set{0,1}$.
Si $f$ n'est pas constante sur $I$, il existe $x,y$ tels que $a<x<y<b$ et $f(x) \neq f(y)$.
On suppose $f(x)=0$ (la démonstration dans le cas $f(x)=1$ est identique).
On considère l'ensemble :
\[ \Gamma = \set{ z \in I ; z \geq x, \forall t \in [x,z], f(t) = 0 } . \]
Cet ensemble est non vide (car il contient $x$) et il est majoré par $y$, il admet donc une borne supérieure que l'on note $c$.
Par continuité de $f$, on a $f(c)=0$. Toujours par continuité, il existe $\veps >0$ tel que :
\[ \forall t \in [c, c+\veps], \quad \delta(f(t), f(c)) \leq \dfrac12 , \]
autrement dit, $f(c+\veps) = 0$ ce qui contredit la définition de $c$ comme borne supérieure de $\Gamma$.
Ainsi, l'hypothèse selon laquelle $f$ n'est pas constante sur $\set{0,1}$ est fausse et donc $I$ est connexe.

\item
Si $I$ est un intervalle quelconque de $\R$, alors on peut toujours trouver un intervalle ouvert $J$ de $\R$ tel que 
$J \subset I \subset \overline{J}$ avec $J$ convexe selon le point (i). Donc $I$ est convexe.
\end{enumerate}
\end{proof}

\begin{coro}[Théorème des valeurs intermédiaires]
Soit $I$ un intervalle de $\R$ et $f$ une fonction continue sur $I$. Alors $f(I)$ est un intervalle.
\end{coro}

\begin{thm}
Soient $E$ un \Rev normé, $[a,b]$ un intervalle de $\R, F : [a,b] \to E$ et $g : [a,b] \to \R$ deux applications continues sur $[a,b]$
et dérivables sur $]a,b[$. Si on a :
\begin{equation} \forall t \in ]a,b[, \quad \norm{F'(t)} \leq g'(t) , \label{eq:1} \end{equation}
alors $\norm{F(b) - F(a)} \leq g(b) - g(a)$.
\end{thm}

\begin{proof}
\begin{enumerate}[label=(\alph*)]

\item
Supposons d'abord que l'inégalité (\ref{eq:1}) est sticte. On peut réécrire cette inégalité comme :
\[ \lim_{\substack{t \to x \\ x < t}} \left ( \norm*{\dfrac{F(t) - F(x)}{t-x}} - \dfrac{g(t)-g(x)}{t-x} \right ) < 0 , \]
donc, par continuité, l'assertion ci-dessous est vraie :
\begin{equation} \forall x \in ]a,b[, \exists y \in ]x,b[ ,\forall t \in [x,y], \quad \norm{F(t) - F(x)} \leq g(t) - g(x) . \label{eq:2} \end{equation}
On va maintenant démontrer que l'assertion suivante est vraie :
\begin{equation} \forall [\alpha, \beta] \subset ]a,b[, \quad \norm{F(\beta) - F(\alpha)} \leq g(\beta) - g(\alpha) , \label{eq:3} \end{equation}
qui permettra d'obtenir le résultat par continuité.
On note :
\[ \Gamma = \set{d \in ]\alpha, \beta] ; \forall t \in [\alpha, d], \norm{F(t) - F(\alpha)} \leq g(t) - g(\alpha) ,  \]
Selon l'assertion (\ref{eq:2}) en prenant $x=\alpha$ on déduit que $\Gamma$ est non vide.
Comme cet ensemble est majoré, il admet une borne supérieure que l'on note $m$. 
On veut simplement démontrer que $m = \beta$. Supposons $m < \beta$.
Par continuité des fonctions $F$ et $g$, on a $\norm{F(m)-F(\alpha)} \leq g(m) - g(\alpha)$, et, d'après l'assertion (\ref{eq:2}) :
\[ \exists y \in ]m,\beta[ ,\forall t \in [m,y], \quad \norm{F(t) - F(m)} \leq g(t) - g(m) , \]
et donc :
\[ \forall t \in [m, y], \quad \norm{F(t)-F(\alpha)} \leq  \norm{F(t) - F(m)} + \norm{F(m)-F(\alpha)} \leq g(t) - g(\alpha) . \]
Autrement dit, $y \in \Gamma$ ce qui est absurde car $y > m = \sup{\Gamma}$.
On a donc $m = \beta$, l'assertion (\ref{eq:3)) est donc vraie sur tout compact inclus dans $]a,b[$, elle est donc vraie (par continuité) 
sur l'intervalle ouvert $]a,b[$, ce qui démontre le résultat dans le cas où l'inégalité (\ref{eq:1}) est sticte

\item
On revient au cas général, prenons $\veps >0$ et définissons la foncton $g_{\veps}$ sur $[a,b]$ par $g_{\veps}(t)=g(t)+\veps(t)$.
On a bien, pour tout $t \in ]a,b[, \norm{F'(t)} < g_{\veps}'(t)$ et donc : 
\[ \norm{F(b) - F(a)} \leq g_{\veps}(b) - g_\veps(a) = g(b) - g(a) + \veps(b-a) , \]
en vertu du point (i). Ceci étant vrai pour tout $\veps >0$, on a bien le résultat attendu.
\end{enumerate}
\end{proof}

\end{document}